
\chapter{Podsumowanie i wnioski}
\section{Wnioski końcowe}
Przedstawione poniżej wnioski opierają się na pomiarach wykonanych na systemie
Linux ogólnego przeznaczenia, który z~założenia nie daje możliwości pracy w twardym systemie wbudowanym czasu rzeczywistego.
Stosowana konfiguracja systemu oraz metodyka pomiarowa pozwoliły
uzyskać rozstęp międzykwartylowy względny poniżej 4\%, co może sugerować
możliwość interpretacji różnic na poziomie 5-7\%. Warto jednak zaznaczyć, że uzyskany rozstęp nie pozwala określić badanego systemu jako takiego, który 
mógłby pracować pod rygorem czasu rzeczywistego.

Najważniejsze wnioski z przeprowadzonych badań prezentują się następująco:

\begin{enumerate}
    \item Największy pojedynczy zysk dało przejście z~\texttt{-O0} na
          \texttt{-O1}. Skala tego zysku zależała od rozmiaru danych - wyniosła
          5,11 przy $N=16384$ i~tylko 2,18 przy $N=1048576$. Różnica ta może
          sugerować, że nawet skrajna różnica w~jakości kodu przestawała mieć
          znaczenie, gdy o~czasie wykonania decydował dostęp do pamięci. Warto
          przy tym wniosku nadmienić, że kompilatory GCC i~Clang przyjmują poziom
          \texttt{-O0} jako domyślną opcję.

    \item Ręczna wektoryzacja z~użyciem instrukcji NEON dała wyraźną przewagę nad
          pozostałymi wariantami radix-2 dla rozmiarów do $N=16384$, ale przy
          $N=1048576$ zysk zanikł całkowicie. Powyżej pojemności pamięci L2
          cztery implementacje radix-2 dzieliło już tylko 2,1\% wobec 65\% przy
          $N=16384$, co oznaczało, że o~wydajności decydowała wtedy organizacja
          dostępu do pamięci, a~nie jakość wygenerowanego kodu. Wyniki te
          pozwoliły uznać jawne użycie instrukcji specyficznych dla architektury
          za skuteczną metodę optymalizacji, wymagającą jednak wiedzy
          wykraczającej poza znajomość samej listy instrukcji. Układ danych
          konieczny do efektywnego wykorzystania rejestrów wektorowych pogarszał
          bowiem lokalność dostępów: wariant NEON notował 5,43\% chybień L1D
          wobec 1,08\% dla \texttt{radix2\_v0}.

    \item Kompilatory GCC i~Clang nie tłumaczyły kodu w~języku C na jednakowy
          program wykonalny. W~wykonanych badaniach Clang generował kod, który
          miał niższą medianę czasów wykonania przy rozmiarach danych wejściowych
          mieszczących się w~pamięci podręcznej, kosztem istotnie większej
          wielkości programu. GCC tracił wektoryzację włączoną wcześniej przez
          \texttt{-O2} po dodaniu pewnej grupy innych flag - dla
          \texttt{radix2\_v1} usuwało ją \texttt{-ffast-math} lub \texttt{-mcpu},
          co wydłużało czas z~0,7562 do~0,8442~ms. Clang zachowywał ją we
          wszystkich zestawach. Sama automatyczna wektoryzacja okazała się
          zjawiskiem rzadkim: wystąpiła w~30 z~80 konfiguracji pod GCC i~w~38 pod
          Clang, a~poziomy \texttt{-O1} oraz \texttt{-Os} nie włączały jej
          w~ogóle. Decydowała o~niej pojedyncza flaga, a~nie poziom optymalizacji.

    \item Kompilator nie naprawiał samodzielnie błędów numerycznych tkwiących
          w~algorytmie. Obserwacja ta była szczególnie widoczna przy zestawieniu
          dokładności dla 16384~próbek. Dla jednej implementacji kompilowanej na
          jedenaście sposobów wartości SNR rozchodziły się najwyżej o~1,31~dB,
          a~zestawy zawierające \texttt{-ffast-math} obniżały wynik o~nie więcej
          niż 0,45~dB. Porównanie między wariantami algorytmu pokazało zupełnie
          inną tendencję. Siedem wariantów osiągnęło 134,89--138,48~dB, natomiast
          \texttt{radix2\_v0} tylko 91,52~dB, czyli o~44~dB mniej, co wskazywało
          na znaczące niedokładności wyników. Żaden poziom optymalizacji ani
          zestaw flag tej różnicy nie zmniejszył, poprawa nastąpiła dopiero po
          zmianie implementacji.

    \item Powyżej poziomu \texttt{-O1} zyski były niesystematyczne. Najlepszy
          zestaw flag skracał czas względem \texttt{-O1} o~kolejne 0,5-29\%, przy
          medianie około 12\%, jednak w~szesnastu badanych konfiguracjach
          najkrótszy czas osiągało osiem różnych zestawów, a~\texttt{-O3} nie
          wygrał w~żadnej z~nich. Podniesienie poziomu optymalizacji potrafiło
          przy tym zaszkodzić - dla \texttt{radix4\_naive} pod GCC poziomy
          \texttt{-O2} i~\texttt{-O3} okazały się o~około 40\% wolniejsze od
          \texttt{-O1}. Uniwersalnie najlepszy zestaw flag zatem nie istniał:
          pewne było jedynie, że \texttt{-O0} dawał najgorszy wynik, a~powyżej
          tego poziomu o~wyniku decydowały konkretna implementacja i~rozmiar
          przetwarzanych danych. Dobór flag był więc zadaniem pomiarowym,
          a~uzyskany wynik należy potwierdzić analizą kodu wynikowego oraz metryk
          optymalizacji.

    \item Rozmiar przetwarzanych danych warunkował dobór schematu algorytmu.
          Kryterium oparte na liczbie operacji arytmetycznych zachowywało ważność
          wyłącznie wtedy, gdy zbiór roboczy mieścił się w~pamięci podręcznej
          pierwszego poziomu - przy $N=1024$ oba bufory zajmowały 16~kB wobec
          32~kB pojemności L1D i~najszybszy był \texttt{splitradix}, schemat
          o~najmniejszej liczbie mnożeń. Przy $N=1048576$ ten sam schemat wypadł
          najgorzej, a~zdecydowanie najszybszy okazał się \texttt{stockham},
          przechodzący przez tablicę liniowo - koszt pojedynczej operacji
          motylkowej wyniósł dla niego 8,13~ns wobec 28,50~ns dla
          \texttt{splitradix}. Jak wywnioskowano z~analizy plików asemblera, zysk
          zaobserwowany dla wariantu Stockhama wynikał prawdopodobnie ze sposobu
          dostępu do pamięci, a~nie, jakby sugerowała intuicja, z~liczby
          wykonywanych działań. Opisywany wariant, jako jedyny z~badanych, nie
          wymagał wstępnej permutacji bitowej, a~więc dodatkowego przebiegu po
          całej szesnastomegabajtowej tablicy. Działając na dwóch buforach,
          w~każdym etapie czytał i~zapisywał kolejne komórki obu tablic, podczas
          gdy warianty działające w~miejscu odwoływały się do znacznie oddalonych
          elementów. Skutkiem tego rozwiązania było prawie dwukrotnie większe
          zapotrzebowanie na pamięć, co tłumaczyło, dlaczego ten sam schemat nie
          okazał się najszybszy przy $N=1024$.

    \item Zrównoleglenie opłacało się dopiero wtedy, gdy praca przypadająca na
          jeden etap przekraczała koszt mechanizmu synchronizacyjnego
          (w~przypadku niniejszej pracy bariery). Przy $N=4096$, gdzie na jedną
          transformatę przypadało dwanaście barier, żaden (wyłączając warianty
          z~flagą \texttt{-O0}) z~wariantów nie osiągnął przyspieszenia. Przy
          $N=65536$ przyspieszenie sięgnęło natomiast 1,94. W~stosowanej
          platformie sprzętowej najniższe czasy wykonania uzyskiwano przy
          stosowaniu czterech wątków, zgodnie z~liczbą rdzeni platformy. Sam
          podział pracy między wątki okazał się równie istotny jak wybór
          schematu: przy niezmienionym algorytmie radix-2 podział płaskiej
          przestrzeni motylków dawał 1,94 wobec 1,51 dla podziału po blokach,
          który zmniejszał się znacząco w~końcowych etapach, gdy liczba bloków
          transformaty spadała poniżej liczby wątków. Schemat najlepszy dla
          wersji jednowątkowej nie musiał być najlepszy po zrównolegleniu -
          \texttt{radix4\_par}, najsłabszy przy $N=65536$ z~wynikiem 1,28, okazał
          się najlepszy przy $N=1048576$ z~wynikiem 1,65, dzięki o~połowę
          mniejszej liczbie etapów, a~więc i~barier. Przy analizie
          zaimplementowanych algorytmów wielowątkowych zauważono, że dobór flag
          dawał mniejsze korzyści. Dla obu wariantów radix-2 rozrzut czasów
          między zestawami nie przekraczał 2,4\%, a~więc pozostawał poniżej progu
          istotności, a~jedynym wyjątkiem był \texttt{radix4\_par} z~zestawem
          \texttt{unroll}, dla którego przewaga wyniosła 8,3\%. Wniosek ten mógł
          być zgodny z~przewidywaniami teoretycznymi, ponieważ stosowane flagi
          kompilacji nie miały bezpośredniej możliwości optymalizacji komunikacji
          międzywątkowej, czyli barier.
\end{enumerate}

Uzyskane wyniki pozwoliły porównać metody na zwiększenie wydajności co do rzędu
wielkości. Dobór schematu obliczeń dopasowanego do rozmiaru danych oraz
odejście od poziomu \texttt{-O0} dały efekty tego samego rzędu - odpowiednio
około 3,5 raza między najlepszym a~najgorszym schematem przy $N=1048576$ oraz
od 2,18 do~5,11 raza przy zmianie poziomu optymalizacji. Warto zaznaczyć przy
tym, że drugi z~nich nie wymagał zmiany ani jednej linii kodu programu
źródłowego. Wyraźnie mniej wniosło strojenie flag powyżej poziomu \texttt{-O1},
warte zwykle około 12\%.

Kompilator można w~takim razie traktować jako realne narzędzie, lecz nie
zastępuje on kluczowych decyzji projektowych. Praktyczny wniosek
z~przeprowadzonych badań jest taki, że optymalizację należy zaczynać od doboru
algorytmu i~odpowiedniej organizacji danych, ponieważ o~czasie działania
programu nie decyduje wyłącznie liczba wykonywanych działań arytmetycznych,
lecz w~porównywalnym stopniu wzorzec dostępu do pamięci.

\section{Kierunek dalszych badań}
Największy potencjał do dalszych badań ma jawne użycie instrukcji specyficznych
dla stosowanej architektury sprzętowej. W~pracy zastosowano je tylko w~jednym wariancie i~przy jednym
układzie danych, co ujawniło konflikt między wektoryzacją a~lokalnością dostępów do pamięci.
Wartościowym kierunkiem badań może być sprawdzenie połączenia instrukcji wektorowych ze schematem
obliczeń o~korzystnym wzorcu dostępu do pamięci, z potencjalnym sprawdzeniem działania proponowanego schematu 
w~wariantach wielowątkowych.

Na uwagę zasługują również badania samego mechanizmu 
synchronizacji wariantów wielowątkowych, zwłaszcza w kotekście stosowania flag optymalizajci. Istotną informacją byłoby w tym przypadku
 wyznaczonenie prógów opłacalności stosowania wariantów wielowątkowych. Przy okazji realizacji proponowanych badań możnaby sprawdzić opłacalność blokowania obliczeń pod
rozmiar pamięci podręcznej, stosowane w~bibliotece FFTW \cite{FFTWPaper}.